\section{Discussion}
\label{sec:discussion}

\subsection{Why Might Self-Editing Hurt Small Models?}

Our experiments measure \emph{that} self-editing lowered success; they were not designed to test
\emph{why}. Below we describe three explanations that fit what we see in the logs. None of them has
been tested with a controlled intervention, so we present them as explanations that are consistent
with our observations, not as established causes.

\paragraph{Wrong lessons.}
The refiner learns from the worker's recent steps. When the worker fails most tasks, these steps
are full of mistakes, and a small refiner can turn the mistakes into rules. For example, after a
few tasks, qwen's own refiner wrote this lesson into its strategy (shortened):
\begin{quote}\small
\emph{``The agent attempted \texttt{cd /task} but the task folder doesn't exist [\ldots]
Lesson: Always use the exact local path provided (\texttt{/task}) and verify it exists [\ldots]
Install dependencies using \texttt{pip install pytest}.''}
\end{quote}
Both pieces of advice are wrong: there is no \texttt{/task} folder (the agent is already in the task
folder), and the sandbox has no network, so \texttt{pip install} always fails. Such rules are then
shown to the worker on every following task. The Reflexion baseline shows the same pattern: its
reflections are confident explanations of failures that are often not the real cause. One
reflection after a failed CSV task claimed the answer was rejected because \texttt{382.35} and
\texttt{382.350} ``do not match''. In fact the checker compares the values as numbers, so the
format does not matter; the real problem was that the total itself was wrong (the correct total is
351.25). This matches the finding that models struggle to correct themselves without a reliable
outside signal~\cite{huang2024cannot}. We did not count how many lessons were wrong across all
runs, so these are examples, not a measured rate.

\paragraph{Longer prompts and more steps.}
Every edit adds text to the prompt. With C1, qwen's prompts became three times longer, and the share
of tasks that ran out of steps rose from 48\% to 88\%. Small models are known to use long contexts
poorly~\cite{liu2024lost}. Across our conditions, the ones that use the most steps tend to have the
lowest success (Reflexion and structured errors), and C8, the only learning condition with about the
same number of steps as the fixed harness, is also the only one with about the same success. This
is a correlation across conditions. Prompt length and prompt content change together in our runs,
so we cannot tell whether the length itself, or the wrong advice inside it, causes the extra steps.
A fixed harness padded to the same length with neutral text would separate the two; we have not
run this control.

\paragraph{New tools create new mistakes.}
Skills and sub-agents are only useful if the worker calls them correctly. Qwen rarely called the
skills it had, and gemma sometimes called a skill's name as if it were a tool. Each new piece of
harness is one more thing a small model can use wrongly.

\subsection{Why Might the Techniques Help Only Together?}

Each technique targets one problem, but C1 may have several problems at once. A better refiner
still writes edits that nobody checks; a gate still lets most edits through when the model fails
almost everything; retrieval may still retrieve wrong lessons; and verified skills are still rarely
called. In C8 these parts may cover for each other: the larger refiner may write fewer wrong rules,
the gate removes some edits, retrieval keeps the prompt shorter, and broken skills are filtered
out. This is one reading of the results; we have not measured each part's contribution inside C8.
The difference between C8 and C1 on qwen is nominally significant and consistent across seeds, but
not significant after correction (Section~\ref{sec:results}).

Two further points matter for reading C8. First, C8 is not self-improvement by the small model: a
larger model writes the edits. Part of its benefit may come from that model's skill rather than from
the safeguards; the single separate-refiner condition (22\%) suggests the larger refiner alone is
not enough, but we cannot say how much it adds inside C8. Second, structured errors alone gave the
lowest success of the single techniques, yet it is part of the best condition. One possible reason
is that it makes every failed tool result longer, which costs more when C1's prompts are already
long. A study that removes one technique at a time from C8 would answer both questions; we leave it
for future work.

\subsection{The Gate Depends on the Model's Strength}

The regression gate behaves very differently on the two models:
\begin{itemize}
  \item For the \textbf{weak} model, it is too \emph{weak}. Qwen fails most tasks, so the old and new
        harness usually both pass 0 or 1 of the 3 rerun tasks. A tie counts as ``not worse'', so most
        edits are kept, whether they help or not.
  \item For the \textbf{strong} model, it is \emph{strict}. Gemma passes almost every task, so a
        single extra failure in the rerun undoes the edit. It undid 5 of 6 edits, and 2 of 3 seeds
        ended with the starting harness.
\end{itemize}
In both cases, three reruns are too few to measure the effect of an edit, and since reruns are not
exact replays (Section~\ref{sec:setup}), the decision partly reflects chance. A gate with more
reruns, or one that rejects ties and needs a clear gain to accept an edit, would be more reliable
but also more expensive; we did not test such variants. For gemma, keeping the harness unchanged was
a reasonable outcome on our task set: the fixed harness is already at 92\%, so there was little to
gain.

\subsection{What This Means in Practice}

For someone building an agent on a small local model, our results suggest:
\begin{enumerate}
  \item \textbf{Do not turn on self-editing by default.} Without safeguards it lowered success for
        both models we tested.
  \item \textbf{If you use it, add supervision.} Test every edit before keeping it, and consider
        letting a larger model write the edits. On our tasks this combination removed the loss, but
        it did not give a clear gain.
  \item \textbf{Keep the prompt short.} Retrieve only the relevant memory, cap every part of the
        harness, and watch the number of steps per task as an early warning sign.
  \item \textbf{Measure against a fixed harness}, on tasks the system did not learn from. A
        self-improving agent is only useful if it beats doing nothing.
\end{enumerate}

\subsection{The Capability Floor}

We defined the floor operationally (Section~\ref{sec:intro}): a model is below it, for a given loop
and task set, if the unmodified loop lowers its success compared with a fixed harness. Both of our
models are below it on our tasks: a 0.8B model at 32\% baseline success and a stronger small model
at 92\%. Two models are two points, not a curve. They show that the floor, for this loop and these
short tasks, is above both models; they do not show where it lies, or whether it depends mainly on
model size, baseline success or task length.

Within this limit, our results suggest what being below the floor looks like. The models do make
edits, and many of them. What they lack, judging from the logs, is a way to tell good edits from bad
ones, so the bad ones are kept. The techniques that helped most add checks from outside the small
model: a larger refiner, a test before keeping an edit, and code that must run. With these checks,
the weak model returned to the fixed-harness level, with a suggestive but unconfirmed sign of better
transfer to new tasks. Without them, self-editing made things worse.
